\section{Implementation}
\label{sec:impl}
% Line counts: `grep -v '^\s*(//|$)' | wc -l` at the commit named in
% tables/numbers.tex; generated files (bpf2go bindings, weights.rs) excluded
% where stated. Module sizes from `ls -l rules/wasm`.

The daemon is about 3,400 lines of Go (plus 1,600 of tests) on
wasmtime-go~v17~\cite{wasmtime} and cilium/ebpf; the kernel side is 190
lines of C compiled with clang and embedded through \texttt{bpf2go}, so the
binary carries its own eBPF objects and the record layout is checked
against the C struct at build time. Rules are Rust compiled to
\texttt{wasm32-wasip1}. A 620-line crate, \texttt{sentinel-abi}, holds the
record layout, the binary encoding and the \texttt{EventFields} trait and is
shared by every rule; \texttt{suspicious\_shell} is 330 lines and
\SI{91}{kB} of Wasm, \texttt{ml\_detector} is 1,900 lines plus a generated
19,000-line weights file and \SI{334}{kB} of Wasm. The ten adversarial
modules of \S\ref{sec:eval:e6} total 220 lines. Nothing in the host is
specific to Rust: any toolchain that can export the five functions of
\S\ref{sec:design:engine} produces a loadable rule.

Three implementation details matter for the measurements. First, the
engine reuses one guest buffer per rule across calls (it asks the guest to
allocate once and grows only when an event is larger), so the steady-state
cost of a call is one copy plus the invocation. Second, fuel and epoch
enforcement are enabled at engine construction; both add a small, fixed
per-instruction and per-tick cost that E2's trivial-guest benchmark
includes. Third, the feature extractor in \texttt{ml\_detector} is a
line-by-line port of the Python function used for training, including its
iteration order over hash maps, and a golden test with 18 recorded inputs
runs in \texttt{cargo test} so that the parity of \S\ref{sec:eval:e2} is
checked on every build. The same source tree builds the replay collector,
the metrics endpoint and the evaluation harnesses; every table in this
paper is regenerated by one \texttt{make} target from the JSON files it
cites.
